\section{The dimension-nine comparison}\label{sec:dimension-nine}

There are two dimension-sensitive steps in the proof.  The classification
of the end profiles used Theorem~\ref{thm:stable} in dimension seven.
The cone argument used the absence of a singular ray in ambient dimension
eight.  In dimension nine, the singular-set bound allows dimension one,
so a vertical cylinder over a singular minimizing cone in $\R^8$
is no longer excluded.  The actual construction of del Pino, Kowalczyk,
and Wei realizes precisely this possibility.

Here the numbering refers to their published paper
\citep{DPKW2011}.  Its Theorem~1, pp.~1488--1489, constructs, for
sufficiently small fixed $\alpha>0$, a bounded entire smooth solution $u_\alpha$ in
$\R^9$, strictly increasing in $x_9$, associated with the nonaffine
entire minimal graph
\[
 \Gamma_\alpha=\{(p,F_\alpha(p)):p\in\R^8\},\qquad
 F_\alpha(p)=\alpha^{-1}F(\alpha p).
\]
It satisfies
\begin{equation}
 |u_\alpha(x)|\longrightarrow1
 \quad\text{as }\dist(x,\Gamma_\alpha)\longrightarrow\infty.
 \label{eq:monotone-dpkw-clearing}
\end{equation}
The graph is a refinement of the Bombieri--De Giorgi--Giusti construction
\citep{BDG1969}.  With
$p=(\xi,\zeta)\in\R^4\times\R^4$ and
$(|\xi|,|\zeta|)=(r\cos\theta,r\sin\theta)$, Lemma~2.1
and Theorem~2 of \citet[pp.~1492--1493]{DPKW2011} give
\begin{equation}
 \begin{gathered}
 F_0(r,\theta)=r^3g_{\rm BDG}(\theta),\qquad
 g_{\rm BDG}(\theta)>0\quad(\pi/4<\theta<\pi/2),\\
 F_0\le F\le F_0+C r^{-\beta}\min\{F_0,1\}
       \quad\text{in this sector for }r>R_0,
       \qquad 0<\beta<1,\\
 F(\zeta,\xi)=-F(\xi,\zeta).
 \end{gathered}
 \label{eq:monotone-dpkw-graph}
\end{equation}
The angular notation $g_{\rm BDG}$ is unrelated to the heteroclinic
$g$.  In arXiv version 0806.3141v2 the corresponding main and graph
theorems are numbered 1.1 and 3.1, respectively.

We spell out the consequences needed for the comparison.  For a fixed
base point $p$, continuity bounds $F_\alpha$ on every bounded
neighborhood of $p$.  Hence
$\dist((p,t),\Gamma_\alpha)\to\infty$ as $|t|\to\infty$:
otherwise a graph point a bounded distance away would have a bounded
base coordinate and an unbounded height.  Bounded strict monotonicity
gives distinct vertical limits, while
\eqref{eq:monotone-dpkw-clearing} puts both limits in $\{-1,1\}$.
They are therefore $-1$ and $1$.  In particular
$-1<u_\alpha<1$ everywhere and each vertical line has a unique zero,
of height $h_\alpha(p)$.  Equation
\eqref{eq:monotone-dpkw-clearing} also implies that all these zeros
lie within a fixed distance $D_\alpha<\infty$ of $\Gamma_\alpha$.

Fix $p$ with $0<|\xi|<|\zeta|$ and fix $t\in\R$.
The angular direction stays in a compact subinterval of the positive
sector for all points $Rp+O(1)$ as $R\to\infty$.  The lower bound
in \eqref{eq:monotone-dpkw-graph}, with $\alpha$ fixed, gives
\begin{equation}
 F_\alpha(Rp+O(1))\ge c_{p,\alpha}R^3
 \quad\text{for all sufficiently large }R.
 \label{eq:monotone-dpkw-cubic}
\end{equation}
A graph point within $D_\alpha+1$ of the zero
$(Rp,h_\alpha(Rp))$ has base coordinate $Rp+O(1)$, so
\eqref{eq:monotone-dpkw-cubic} implies
$h_\alpha(Rp)/R\to+\infty$.  Thus $u_\alpha(Rp,Rt)<0$ for
large $R$.  Moreover the distance from $(Rp,Rt)$ to the graph tends
to infinity: a bounded-distance subsequence would contradict
\eqref{eq:monotone-dpkw-cubic}.  By
\eqref{eq:monotone-dpkw-clearing},
\[
 u_\alpha(Rp,Rt)\longrightarrow-1
           \quad\text{if }0<|\xi|<|\zeta|.
\]
The oddness of $F$ and the same argument give limit $+1$ in the
opposite sector $0<|\zeta|<|\xi|$.  The excluded equality set and
coordinate axes have Lebesgue measure zero.  Bounded convergence now
shows local $L^1$ convergence of the rescaled solutions to the two-phase
function with interface
\begin{equation}
 \mathcal C_9=
      \{(\xi,\zeta):|\xi|=|\zeta|\}\times\R.
 \label{eq:monotone-dpkw-cylinder}
\end{equation}
This cone has the entire vertical axis as its singular set and is
invariant under vertical translations.  Its directional normal
component is zero on its regular part.  Thus it falls into exactly the
case that was ruled out by the singular-set bound in
Lemma~\ref{lem:monotone-flat-cone} in dimension eight.

The same solutions are minimizers by
Lemma~\ref{lem:monotone-minimizer}, applied to their constant end
profiles, and consequently have the corresponding $O(R^8)$ energy
growth.  Thus neither the derived minimizing property nor this energy
bound conflicts with the construction.  Their nonplanar phase limit
also precludes one-dimensionality.  The construction extends constantly
in additional coordinates to every $N>9$, as explained in
\citet[p.~1490]{DPKW2011}: the added second derivatives vanish, and
strict monotonicity in $x_9$ is preserved.  It supplies no
eight-dimensional example.
